\documentclass[11pt,a4paper]{article}

\usepackage[margin=25mm]{geometry}
\usepackage{amsmath,amssymb}
\usepackage{array,longtable}
\usepackage{graphicx}
\usepackage[hidelinks]{hyperref}
\hypersetup{
  pdftitle={Adaptive Signal Acquisition: Living Technical Report},
  pdfauthor={Adaptive Signal Acquisition project}
}

\setlength{\parindent}{0pt}
\setlength{\parskip}{0.55em}
\raggedbottom
\emergencystretch=1em
\setcounter{secnumdepth}{3}
\setcounter{tocdepth}{2}
\renewcommand{\arraystretch}{1.18}

\newcommand{\R}{\mathbb{R}}
\newcommand{\E}{\mathbb{E}}
\newcommand{\Var}{\operatorname{Var}}
\newcommand{\Cov}{\operatorname{Cov}}
\newcommand{\supp}{\operatorname{supp}}
\newcommand{\tr}{\operatorname{tr}}
\newcommand{\argmin}{\operatorname*{arg\,min}}
\newcommand{\argmax}{\operatorname*{arg\,max}}
\newcommand{\norm}[1]{\left\lVert #1\right\rVert}
\newcommand{\trans}{^{\mathsf T}}

\title{\textbf{Adaptive Signal Acquisition}\\[0.4em]
\large A Living Technical Report on Sparse Reconstruction and Sequential Sensing}
\author{Project technical record}
\date{Version 3 --- Stages 1 through 4B --- 27 September 2026}

\begin{document}
\maketitle

\begin{abstract}
This report develops a reproducible simulation study of sparse-signal acquisition under a
limited noisy-measurement budget.  The central question is whether a sequential sensing
policy, which chooses each new measurement after seeing earlier observations, improves
reconstruction over fair non-adaptive baselines.  Stages 1 and 2 establish the deterministic
simulator, LASSO reconstruction, and fixed and random references.  Stage 3 introduces a
Gaussian information score augmented by a provisional-LASSO magnitude proxy.  Stage 4A
shows that the proxy is harmful: it rewards alignment with the current estimate rather than
uncertainty.  Stage 4B replaces it with an observation-dependent bootstrap covariance of
plausible sparse reconstructions.  In 100 paired trials the new method removes the large
failure of the first adaptive policy and is statistically competitive with random and
Gaussian-only sensing, but it does not yet provide a significant improvement.  The report
states the mathematical model, derives each score and update, defines every experimental
quantity, records negative results, and identifies the next controlled experiments.
\end{abstract}

\tableofcontents
\newpage

\section{Research question and scope}

The primary question is:

\begin{quote}
When only a limited number of noisy linear measurements may be taken, can selecting each
new measurement using the previous action--observation history reconstruct a sparse signal
more accurately than selecting all measurements in advance?
\end{quote}

The operational endpoint is normalized mean squared error (NMSE).  A comparison is fair
only when the methods share the hidden-signal distribution, candidate measurement
dictionary, noise law, measurement budget, final reconstruction algorithm, and paired
random draws.  The project currently studies a synthetic, directly sparse vector.  It is a
controlled algorithmic test bed, not yet a claim about a particular physical sensor.

\section{Notation and mathematical model}

\subsection{Notation}

\begin{longtable}{>{\raggedright\arraybackslash}p{0.19\textwidth}p{0.73\textwidth}}
\hline
Symbol & Definition \\
\hline
$N$ & Dimension of the unknown signal; $N=64$ in the primary experiments. \\
$k$ & Number of non-zero coordinates; $k=4$ in the primary experiments. \\
$x\in\R^N$ & Hidden signal.  It is available to the simulator but never to a sensing policy. \\
$\supp(x)$ & Support of $x$, namely the indices of its non-zero entries. \\
$D\in\R^{M\times N}$ & Candidate measurement dictionary; $M=128$ normalized rows. \\
$a_t\in\R^N$ & Column-form measurement vector selected at sequential step $t$; $a_t\trans$ is a row of $D$. \\
$B$ & Measurement budget, or number of selected rows. \\
$A_t\in\R^{t\times N}$ & Sensing matrix formed by stacking the first $t$ selected rows. \\
$y_t\in\R$ & Scalar noisy observation returned for action $a_t$. \\
$y_{1:t}\in\R^t$ & Vector containing the first $t$ observations. \\
$\epsilon_t$ & Independent Gaussian measurement noise with variance $\sigma^2$. \\
$\hat{x}_t$ & Provisional sparse reconstruction after $t$ observations. \\
$\mu_t,\Sigma_t$ & Mean and covariance of the Gaussian uncertainty approximation. \\
$\lambda$ & LASSO sparsity-penalty coefficient; it is unrelated to the sensing weight $\gamma$. \\
$\gamma$ & Stage 3/4A weight on the provisional-LASSO magnitude proxy. \\
$R$ & Number of bootstrap reconstructions in Stage 4B; $R=12$. \\
$\rho$ & Stage 4B Gaussian shrinkage weight; $\rho=0.25$. \\
\hline
\end{longtable}

\subsection{Sparse hidden signal}

A vector is $k$-sparse when exactly $k$ coordinates are non-zero:
\begin{equation}
x\in\R^N,
\qquad \norm{x}_0 = |\supp(x)| = k \ll N .
\label{eq:sparsity}
\end{equation}
The notation $\norm{x}_0$ counts non-zero entries; it is not a norm in the strict algebraic
sense.  In each trial, the simulator samples $k$ support indices uniformly without
replacement and assigns independent amplitudes from $\mathcal{N}(0,1)$.  More generally a
signal may be sparse in a known basis $\Psi$, so that $x=\Psi z$ and $z$ is sparse.  The
current implementation uses $\Psi=I$, hence $z=x$.

\subsection{Candidate actions and noisy measurements}

Each candidate row contains $N$ independent Rademacher entries ($-1$ or $+1$ with equal
probability) and is normalized to unit Euclidean length.  Thus
\begin{equation}
D=\begin{bmatrix}d_1\trans\\ \vdots\\ d_M\trans\end{bmatrix}\in\R^{M\times N},
\qquad \norm{d_j}_2=1,
\qquad M=128 .
\label{eq:dictionary}
\end{equation}
The candidate count $M$ is not the signal dimension $N$ and is not the budget $B$.  It is
the number of allowed experiments from which a policy chooses.  If indices
$i_1,\ldots,i_B$ are selected, the sensing matrix is
\begin{equation}
A_B=\begin{bmatrix}a_1\trans\\a_2\trans\\\vdots\\a_B\trans\end{bmatrix}
=\begin{bmatrix}d_{i_1}\trans\\d_{i_2}\trans\\\vdots\\d_{i_B}\trans\end{bmatrix}
\in\R^{B\times N} .
\label{eq:sensing-matrix}
\end{equation}
At step $t$, the environment returns one scalar:
\begin{equation}
y_t=a_t\trans x+\epsilon_t,
\qquad \epsilon_t\stackrel{\text{i.i.d.}}{\sim}\mathcal{N}(0,\sigma^2),
\qquad \sigma=0.10 .
\label{eq:scalar-measurement}
\end{equation}
Stacking the observations gives the standard linear inverse problem
\begin{equation}
y_{1:B}=A_Bx+\epsilon .
\label{eq:linear-model}
\end{equation}
A measurement is therefore a weighted sum of all coordinates, not necessarily a direct
sample of one coordinate.  The simulation represents what a physical linear measurement
device would do after calibration supplies its response coefficients.

\section{Stage 1: deterministic foundations}

Stage 1 creates $x$, $D$, a fixed action subset, $A$, clean measurements $Ax$, a Gaussian
noise vector, and observations $y$.  A seeded NumPy generator makes the arrays exactly
reproducible.  The main shapes are
\begin{equation}
x:(64,),\qquad D:(128,64),\qquad A:(B,64),\qquad y:(B,).
\end{equation}
No reconstruction or adaptivity is used at this stage.  Its purpose is to make every array
in Equation~\eqref{eq:linear-model} inspectable and to ensure that later differences come
from algorithms rather than accidental changes in random data.

\IfFileExists{artifacts/toy_run/toy_run.png}{
\begin{figure}[htbp]
\centering
\includegraphics[width=0.94\textwidth]{artifacts/toy_run/toy_run.png}
\caption{Stage 1 visual check of the sparse signal, clean measurements, noise, and observations.}
\end{figure}}{}

\section{Stage 2: sparse reconstruction and non-adaptive baselines}

\subsection{LASSO objective}

After $B$ measurements the project reconstructs the signal with LASSO:
\begin{equation}
\hat{x}
=\argmin_{z\in\R^N}
\left\{
\frac{1}{2}\norm{A_Bz-y_{1:B}}_2^2+\lambda\norm{z}_1
\right\} .
\label{eq:lasso}
\end{equation}
Here $z$ is a variable representing a candidate reconstruction; it is not the observed
vector $y$.  The first term is small when $z$ explains the data.  The second term,
$\norm{z}_1=\sum_{j=1}^N|z_j|$, discourages many non-zero coefficients.  The primary final
solver uses $\lambda=0.08$, at most 5000 iterations, and tolerance $10^{-7}$.

\subsection{ISTA derivation}

Let $f(z)=\tfrac12\norm{A_Bz-y}_2^2$.  Its gradient is
\begin{equation}
\nabla f(z)=A_B\trans(A_Bz-y).
\label{eq:lasso-gradient}
\end{equation}
The gradient is Lipschitz with constant $L=\norm{A_B}_2^2$, where $\norm{A_B}_2$ is the
largest singular value.  Iterative shrinkage-thresholding (ISTA) performs
\begin{align}
u^{(q)} &= z^{(q)}-\frac{1}{L}\nabla f\!\left(z^{(q)}\right), \\
z^{(q+1)} &= \mathcal{S}_{\lambda/L}\!\left(u^{(q)}\right),
\label{eq:ista}
\end{align}
where soft thresholding acts coordinatewise:
\begin{equation}
\mathcal{S}_{\theta}(u_j)
=\operatorname{sign}(u_j)\max\{|u_j|-\theta,0\} .
\label{eq:soft-threshold}
\end{equation}
The gradient step improves data agreement; soft thresholding is the proximal step for the
$\ell_1$ term and sets sufficiently weak coordinates to zero.

\subsection{Fixed and random reference policies}

The \emph{fixed} policy uses the same precomputed row order in every trial.  The
\emph{random} policy draws a row order before observing any $y_t$ values.  Both are
non-adaptive.  Within a trial all methods receive the same hidden $x$ and the same noise
prefix; only their selected rows differ.  A budget $B=16$ means 16 total measurements, not
four groups of 16 and not 16 observations repeated across trials.

\subsection{Metrics}

The principal metric is
\begin{equation}
\operatorname{NMSE}(x,\hat{x})
=\frac{\norm{x-\hat{x}}_2^2}{\norm{x}_2^2+10^{-12}} .
\label{eq:nmse}
\end{equation}
The $10^{-12}$ term is only a numerical guard.  Lower NMSE is better.  Support recovery is
reported using precision, recall, and their harmonic mean:
\begin{equation}
F_1=\frac{2PR}{P+R},\qquad
P=\frac{|\widehat{S}\cap S|}{|\widehat{S}|},\qquad
R=\frac{|\widehat{S}\cap S|}{|S|},
\label{eq:f1}
\end{equation}
where $S=\supp(x)$ and $\widehat{S}=\{j:|\hat{x}_j|>10^{-3}\}$.

\IfFileExists{artifacts/stage2_baselines/nmse_vs_budget.png}{
\begin{figure}[htbp]
\centering
\includegraphics[width=0.88\textwidth]{artifacts/stage2_baselines/nmse_vs_budget.png}
\caption{Stage 2 baseline behavior.  Reconstruction error decreases as the budget grows.}
\end{figure}}{}

\section{Stage 3: first sequential information-guided policy}

\subsection{History and causal selection}

At time $t$, the available history is
\begin{equation}
h_t=(A_t,y_{1:t}) .
\label{eq:history}
\end{equation}
A valid policy selects $a_{t+1}=\pi(h_t)$ before the simulator evaluates
$a_{t+1}\trans x+\epsilon_{t+1}$.  The selector has no hidden-signal argument.  This ordering
prevents truth leakage.

\subsection{Gaussian information score}

Assume temporarily that the conditional belief is Gaussian:
\begin{equation}
x\mid h_t\sim\mathcal{N}(\mu_t,\Sigma_t).
\label{eq:gaussian-belief}
\end{equation}
For a candidate action $a$, the noiseless scalar $a\trans x$ has variance
\begin{equation}
\Var(a\trans x\mid h_t)
=\E\!\left[(a\trans(x-\mu_t))^2\mid h_t\right]
=a\trans\Sigma_t a .
\label{eq:projected-variance}
\end{equation}
Because independent noise adds variance, $y\mid h_t,a$ has variance
$a\trans\Sigma_ta+\sigma^2$.  The entropy of a scalar Gaussian with variance $v$ is
$\tfrac12\log(2\pi e v)$.  Hence the mutual information is
\begin{align}
I(x;y\mid h_t,a)
&=H(y\mid h_t,a)-H(y\mid x,h_t,a)\\
&=\frac12\log\!\left(1+\frac{a\trans\Sigma_ta}{\sigma^2}\right).
\label{eq:mutual-information}
\end{align}
The implementation omits the positive factor $\tfrac12$ because it cannot change the
ranking.  It selects the unused row maximizing
\begin{equation}
s_G(a)=\log\!\left(1+\frac{a\trans\Sigma_ta}{\sigma^2}\right).
\label{eq:gaussian-score}
\end{equation}

\subsection{Gaussian conditioning update}

After observing $(a_t,y_t)$, the exact linear-Gaussian update is
\begin{align}
K_t&=\frac{\Sigma_ta_t}{a_t\trans\Sigma_ta_t+\sigma^2}, \\
\mu_{t+1}&=\mu_t+K_t(y_t-a_t\trans\mu_t), \\
\Sigma_{t+1}&=\Sigma_t-K_ta_t\trans\Sigma_t .
\label{eq:gaussian-update}
\end{align}
$K_t$ is the gain.  The covariance update depends on the selected direction but not on the
numerical observation.  Therefore Gaussian-only sensing is sequential but, with fixed
dictionary and tie breaking, its row sequence is not observation-value adaptive.

\subsection{The adaptive\_v1 magnitude proxy}

To make Stage 3 react to $y$, a provisional LASSO estimate $\hat{x}_t$ was inserted through
\begin{equation}
\widetilde{\Sigma}_t
=\Sigma_t+\gamma\hat{x}_t\hat{x}_t\trans .
\label{eq:v1-covariance}
\end{equation}
For a candidate $a$,
\begin{equation}
a\trans\widetilde{\Sigma}_ta
=\underbrace{a\trans\Sigma_ta}_{g_t(a)}
+\underbrace{\gamma(a\trans\hat{x}_t)^2}_{\ell_t(a)} .
\label{eq:v1-components}
\end{equation}
Thus $\gamma$ controls only the sensing score; it does not change the LASSO penalty
$\lambda$.  The selected action maximizes
\begin{equation}
s_{v1}(a)=\log\!\left(1+\frac{g_t(a)+\ell_t(a)}{\sigma^2}\right).
\label{eq:v1-score}
\end{equation}

\section{Stage 3 result and Stage 4A diagnosis}

The Stage 3 policy was worse than the non-adaptive references at every tested budget.
Table~\ref{tab:stage3} reports mean NMSE over 100 trials.

\begin{table}[htbp]
\centering
\caption{Stage 3 mean NMSE; lower is better.}
\label{tab:stage3}
\begin{tabular}{rrrr}
\hline
$B$ & Fixed & Random & adaptive\_v1 \\
\hline
8  & 0.8324 & 0.8238 & 0.9006 \\
12 & 0.5616 & 0.6010 & 0.7795 \\
16 & 0.4045 & 0.4185 & 0.6499 \\
20 & 0.3024 & 0.3396 & 0.5078 \\
24 & 0.2400 & 0.2792 & 0.4028 \\
32 & 0.1815 & 0.2004 & 0.2724 \\
\hline
\end{tabular}
\end{table}

Stage 4A varied only $\gamma\in\{0,0.1,0.5,1\}$ and recorded the separate components in
Equation~\eqref{eq:v1-components}, row correlation, singular values, condition number, and
provisional reconstruction quality.  The complete 100-trial result is shown in
Table~\ref{tab:stage4a}.

\begin{table}[htbp]
\centering
\caption{Stage 4A mean NMSE for the magnitude-proxy ablation.}
\label{tab:stage4a}
\begin{tabular}{rrrrrr}
\hline
$B$ & Random & $\gamma=0$ & $\gamma=0.1$ & $\gamma=0.5$ & $\gamma=1$ \\
\hline
8  & 0.8263 & 0.7685 & 0.8197 & 0.8643 & 0.8441 \\
12 & 0.5876 & 0.5649 & 0.6188 & 0.7497 & 0.7033 \\
16 & 0.3750 & 0.3788 & 0.4489 & 0.6126 & 0.5589 \\
20 & 0.2843 & 0.2607 & 0.3310 & 0.4868 & 0.4528 \\
24 & 0.2251 & 0.2070 & 0.2473 & 0.3988 & 0.3669 \\
32 & 0.1663 & 0.1443 & 0.1633 & 0.2681 & 0.2312 \\
\hline
\end{tabular}
\end{table}

At $B=32$, Gaussian-only ($\gamma=0$) had paired difference $-0.0220$ from random with
95\% bootstrap interval $[-0.0424,-0.0011]$.  The original $\gamma=1$ policy had difference
$+0.0649$ with interval $[+0.0277,+0.1007]$.  Increasing $\gamma$ raised the mean maximum
row correlation from 0.17 to 0.22 and the mean condition number from 1.79 to 2.52.

The diagnosis follows directly from Equation~\eqref{eq:v1-components}:
$\hat{x}_t\hat{x}_t\trans$ represents estimated magnitude, not uncertainty.  It rewards
rows aligned with what the provisional estimate already believes and can repeatedly measure
similar directions.  A valid replacement should ask where plausible reconstructions
\emph{disagree}.

\IfFileExists{artifacts/stage4a_diagnostics/nmse_gamma_ablation.png}{
\begin{figure}[htbp]
\centering
\includegraphics[width=0.9\textwidth]{artifacts/stage4a_diagnostics/nmse_gamma_ablation.png}
\caption{Stage 4A: larger magnitude-proxy weights degrade reconstruction.}
\end{figure}}{}

\section{Stage 4B: adaptive\_v2 bootstrap uncertainty}

\subsection{Design requirements}

The replacement was implemented as a separate policy; adaptive\_v1 remains unchanged for
comparison.  The new policy must (i) depend on observed values, (ii) represent uncertainty
rather than magnitude, (iii) retain geometric diversity, (iv) never receive the hidden
signal, and (v) be exactly reproducible under a fixed policy seed.

\subsection{Gaussian-diverse burn-in}

Bootstrap uncertainty cannot be estimated before observations exist.  For the first
$b_0=8$ actions, adaptive\_v2 therefore uses Equation~\eqref{eq:gaussian-score}.  This
burn-in collects diverse rows and produces an initial history $(A_{b_0},y_{1:b_0})$.  It is
not random: it greedily covers directions that remain uncertain under $\Sigma_t$.

\subsection{Parametric bootstrap ensemble}

At a post-burn-in step $t$, generate $R=12$ perturbed observation vectors
\begin{equation}
y_{1:t}^{(r)}=y_{1:t}+\eta^{(r)},
\qquad
\eta^{(r)}\sim\mathcal{N}(0,\tau^2\sigma^2 I_t),
\qquad r=1,\ldots,R,
\label{eq:bootstrap-observations}
\end{equation}
with perturbation multiplier $\tau=1$.  For each perturbation, solve
\begin{equation}
\hat{x}_t^{(r)}
=\argmin_z
\left\{
\frac12\norm{A_tz-y_{1:t}^{(r)}}_2^2+\lambda\norm{z}_1
\right\} .
\label{eq:bootstrap-lasso}
\end{equation}
The implementation evaluates all $R$ ISTA problems in one vectorized matrix update.  The
policy-internal solver uses at most 500 iterations and tolerance $10^{-5}$; the final
evaluation still uses the unchanged 5000-iteration LASSO.

Define the ensemble mean and empirical covariance by
\begin{align}
\bar{x}_t&=\frac1R\sum_{r=1}^R\hat{x}_t^{(r)}, \\
\widehat{\Sigma}_{B,t}
&=\frac{1}{R-1}\sum_{r=1}^R
(\hat{x}_t^{(r)}-\bar{x}_t)(\hat{x}_t^{(r)}-\bar{x}_t)\trans .
\label{eq:bootstrap-covariance}
\end{align}
For any candidate $a$,
\begin{equation}
a\trans\widehat{\Sigma}_{B,t}a
=\frac{1}{R-1}\sum_{r=1}^R
\left[a\trans(\hat{x}_t^{(r)}-\bar{x}_t)\right]^2 .
\label{eq:bootstrap-projection}
\end{equation}
Equation~\eqref{eq:bootstrap-projection} is the sample variance of the measurement predicted
by the ensemble.  A large value means plausible sparse explanations of the existing data
predict different outcomes for action $a$.  That is the intended notion of informativeness.

\subsection{Trace matching and shrinkage}

The Gaussian and bootstrap covariances can have very different total scales.  Directly
adding them would make their relative influence depend on solver scale rather than an
explicit design parameter.  The implementation trace-matches the bootstrap covariance:
\begin{equation}
\widehat{\Sigma}_{B,t}^{\mathrm{scaled}}
=\frac{\tr(\Sigma_t)}
{\max\{\tr(\widehat{\Sigma}_{B,t}),\varepsilon_{\mathrm{num}}\}}
\widehat{\Sigma}_{B,t},
\label{eq:trace-matching}
\end{equation}
where $\varepsilon_{\mathrm{num}}$ is machine-scale protection against division by zero.  If
the empirical covariance is numerically zero, the policy falls back to Gaussian-only
sensing.  Otherwise it forms
\begin{equation}
\Sigma_{v2,t}
=\rho\Sigma_t+(1-\rho)\widehat{\Sigma}_{B,t}^{\mathrm{scaled}},
\qquad \rho=0.25 .
\label{eq:v2-shrinkage}
\end{equation}
The Gaussian term supplies a full-rank diversity reference.  This matters because an
$R$-member sample covariance has rank at most $R-1=11$, much smaller than $N=64$.

\subsection{Action score and selection}

Every unused candidate is scored by
\begin{equation}
s_{v2}(a)
=\log\!\left(1+\frac{a\trans\Sigma_{v2,t}a}{\sigma^2}\right),
\qquad
a_{t+1}=\argmax_{a\in\mathcal{D}\setminus\mathcal{U}_t}s_{v2}(a),
\label{eq:v2-score}
\end{equation}
where $\mathcal{D}$ is the candidate set and $\mathcal{U}_t$ contains already used actions.
The selected raw variance is recorded as separate Gaussian and bootstrap contributions:
\begin{equation}
\rho\,a\trans\Sigma_ta
+(1-\rho)\,a\trans\widehat{\Sigma}_{B,t}^{\mathrm{scaled}}a .
\label{eq:v2-components}
\end{equation}
Unlike the Gaussian covariance, $\widehat{\Sigma}_{B,t}$ is a function of $y_{1:t}$.  LASSO
is nonlinear because soft thresholding can change active sets, so changing the observations
can change both the ensemble support patterns and covariance.  A unit test explicitly
verifies that two observation histories produce different bootstrap covariances while using
the same perturbation draws.

\subsection{End-to-end algorithm}

For one trial, adaptive\_v2 performs the following causal loop:
\begin{enumerate}
\item Initialize $\mu_0=0$ and $\Sigma_0=(4/64)I=0.0625I$.
\item For $t<8$, score unused rows with Gaussian variance only.
\item For $t\geq8$, construct the bootstrap ensemble from $(A_t,y_{1:t})$, compute
      Equations~\eqref{eq:bootstrap-covariance}--\eqref{eq:v2-shrinkage}, and score unused rows.
\item Select the maximum-score row.  Only now does the simulator evaluate
      $y_{t+1}=a_{t+1}\trans x+\epsilon_{t+1}$.
\item Append the row and observation, update the Gaussian belief, and repeat.
\item At each requested budget, reconstruct with the common final LASSO and compute metrics.
\end{enumerate}

\section{Stage 4B experiment and result}

\subsection{Frozen protocol}

The study used 100 paired trials, $N=64$, $k=4$, $M=128$, $\sigma=0.10$, budgets
$B\in\{8,12,16,20,24,32\}$, final $\lambda=0.08$, and seed 20260927.  Each method was run
once to $B=32$ per trial; smaller budgets reuse exact prefixes.  Five policies were compared:
fixed, random, Gaussian-only, frozen adaptive\_v1 with $\gamma=1$, and adaptive\_v2.
Confidence intervals are deterministic percentile-bootstrap intervals from 2000 resamples of
the paired trial-level NMSE differences.

\begin{table}[htbp]
\centering
\caption{Stage 4B mean NMSE over 100 paired trials.}
\label{tab:stage4b}
\begin{tabular}{rrrrrr}
\hline
$B$ & Fixed & Random & Gaussian-only & adaptive\_v1 & adaptive\_v2 \\
\hline
8  & 0.7956 & 0.7983 & 0.8210 & 0.8815 & 0.8210 \\
12 & 0.5959 & 0.5243 & 0.5547 & 0.6967 & 0.5627 \\
16 & 0.3894 & 0.3775 & 0.3437 & 0.5759 & 0.3807 \\
20 & 0.2904 & 0.2807 & 0.2707 & 0.4643 & 0.2748 \\
24 & 0.2289 & 0.2281 & 0.2231 & 0.4005 & 0.2228 \\
32 & 0.1541 & 0.1657 & 0.1480 & 0.2788 & 0.1586 \\
\hline
\end{tabular}
\end{table}

\begin{table}[htbp]
\centering
\caption{Paired adaptive\_v2 NMSE difference from random; negative favors adaptive\_v2.}
\label{tab:v2-ci}
\begin{tabular}{rrrr}
\hline
$B$ & Mean difference & 95\% lower & 95\% upper \\
\hline
8  & +0.0227 & -0.0469 & +0.0940 \\
12 & +0.0384 & -0.0312 & +0.1111 \\
16 & +0.0033 & -0.0534 & +0.0651 \\
20 & -0.0059 & -0.0513 & +0.0388 \\
24 & -0.0053 & -0.0469 & +0.0410 \\
32 & -0.0071 & -0.0383 & +0.0332 \\
\hline
\end{tabular}
\end{table}

\IfFileExists{artifacts/stage4b_adaptive_v2/nmse_comparison.png}{
\begin{figure}[htbp]
\centering
\includegraphics[width=0.92\textwidth]{artifacts/stage4b_adaptive_v2/nmse_comparison.png}
\caption{Stage 4B comparison.  adaptive\_v2 removes the large adaptive\_v1 degradation but
does not separate significantly from random or Gaussian-only sensing.}
\end{figure}}{}

\clearpage
\subsection{Interpretation}

The result supports three precise statements.
\begin{enumerate}
\item \textbf{The replacement fixed the identified failure mode.}  At $B=32$, adaptive\_v1
      is worse than random by $+0.1132$ with interval $[+0.0745,+0.1534]$, whereas
      adaptive\_v2 differs by $-0.0071$ with interval $[-0.0383,+0.0332]$.
\item \textbf{Observation-dependent uncertainty did not yet produce a demonstrated gain.}
      Every adaptive\_v2 interval in Table~\ref{tab:v2-ci} crosses zero.  It is therefore
      accurate to call the method competitive, not superior.
\item \textbf{The tested dense random dictionary leaves limited room for adaptivity.}
      Fixed, random, Gaussian-only, and adaptive\_v2 become close at moderate and large
      budgets.  Dense normalized rows already provide broad coverage, so selecting among
      them may matter less than it would for structured dictionaries or clustered supports.
\end{enumerate}

Across the 2400 post-burn-in decisions, the bootstrap term supplied a mean 0.8705 fraction
of the selected raw variance, the empirical covariance had mean rank 10.98 (consistent with
the theoretical maximum 11), mean maximum row correlation was 0.25, mean condition number
was 2.83, and mean post-observation provisional NMSE was 0.3389.  These diagnostics show that the bootstrap
term is active; the null result is not caused by silently falling back to Gaussian-only
sensing.

\IfFileExists{artifacts/stage4b_adaptive_v2/paired_difference_vs_random.png}{
\begin{figure}[htbp]
\centering
\includegraphics[width=0.9\textwidth]{artifacts/stage4b_adaptive_v2/paired_difference_vs_random.png}
\caption{Paired NMSE differences from random sensing with 95\% bootstrap intervals.}
\end{figure}}{}

\IfFileExists{artifacts/stage4b_adaptive_v2/adaptive_v2_diagnostics.png}{
\begin{figure}[htbp]
\centering
\includegraphics[width=0.92\textwidth]{artifacts/stage4b_adaptive_v2/adaptive_v2_diagnostics.png}
\caption{Decision diagnostics.  The dashed vertical line separates the eight-step burn-in
from bootstrap-guided sensing.}
\end{figure}}{}

\section{What should happen next}

Stage 4B is complete as an implementation and base-setting evaluation.  It should not yet be
followed by an LLM orchestrator, because the scientific action space is not sufficiently
understood.  The next work should be a controlled Stage 4C sensitivity study:
\begin{enumerate}
\item vary $\rho\in\{0,0.25,0.5,0.75,1\}$ to measure the diversity--bootstrap trade-off;
\item vary burn-in length and ensemble size $R$;
\item compare observation perturbation with residual and wild bootstrap variants;
\item test lower SNR, different sparsity levels, clustered supports, and structured
      dictionaries in which adaptive selection has a clearer opportunity to help;
\item pre-register a primary setting and confirm selected conclusions with 500 trials.
\end{enumerate}
Only after this policy family and experimental registry are stable should an LLM manage
experiment selection, summarize evidence, or propose bounded follow-up configurations.  The
LLM should not replace the numerical policy, the reconstructor, or the deterministic result
registry.

\section{Implementation map and reproducibility}

\begin{center}
\begin{tabular}{p{0.39\textwidth}p{0.53\textwidth}}
\hline
Path & Responsibility \\
\hline
\texttt{signals.py} & Sparse-signal generation. \\
\texttt{measurements.py} & Candidate dictionary and $y=Ax+\epsilon$. \\
\texttt{reconstruction.py} & LASSO objective and ISTA solver. \\
\texttt{baselines.py} & Fixed and random action orders. \\
\texttt{sequential.py} & Gaussian belief and frozen adaptive\_v1. \\
\texttt{stage4.py} & Stage 4A gamma ablation and diagnostic records. \\
\texttt{adaptive\_v2.py} & Bootstrap ensemble, covariance, shrinkage, and sequential loop. \\
\texttt{stage4b.py} & Paired Stage 4B experiment and bootstrap confidence intervals. \\
\texttt{run\_stage4b\_study.py} & CSV, metadata, and figure generation. \\
\texttt{tests/test\_stage4b.py} & Determinism, causality, observation dependence, and paired-study tests. \\
\hline
\end{tabular}
\end{center}

Run the verification and primary study from the repository root:
\begin{verbatim}
.venv\Scripts\python.exe -m pytest -q
.venv\Scripts\python.exe scripts\run_stage4b_study.py
.venv\Scripts\python.exe scripts\build_living_report.py
\end{verbatim}

The Stage 4B output directory is
\texttt{artifacts/stage4b\_adaptive\_v2/}.  It contains trial-level records, paired summaries,
decision traces, exact metadata, and figures.  Random seeds are separated for dictionary,
fixed order, signal, noise, random baseline, bootstrap policy, and confidence-interval
resampling.  Runtime naturally varies; selected actions and quality metrics reproduce.

\section{Limitations}

The current evidence is limited to directly sparse Gaussian-amplitude signals, a dense
Rademacher candidate dictionary, one noise level, and LASSO reconstruction.  The bootstrap
ensemble perturbs observations conditionally on one measured history and does not represent
all forms of model uncertainty.  Trace matching is an explicit heuristic, not a Bayesian
identity.  The final numerical conclusions apply to the frozen configurations reported here;
they do not establish that adaptive sensing is generally useless or generally superior.

\small
\begin{thebibliography}{9}
\setlength{\itemsep}{0.15em}
\bibitem{tibshirani1996}
R. Tibshirani, ``Regression Shrinkage and Selection via the Lasso,''
\emph{Journal of the Royal Statistical Society: Series B}, 1996.

\bibitem{daubechies2004}
I. Daubechies, M. Defrise, and C. De Mol, ``An Iterative Thresholding Algorithm for Linear
Inverse Problems with a Sparsity Constraint,'' \emph{Communications on Pure and Applied
Mathematics}, 2004.

\bibitem{candes2006}
E. Cand\`es, J. Romberg, and T. Tao, ``Robust Uncertainty Principles: Exact Signal
Reconstruction from Highly Incomplete Frequency Information,''
\emph{IEEE Transactions on Information Theory}, 2006.

\bibitem{efron1993}
B. Efron and R. Tibshirani, \emph{An Introduction to the Bootstrap}, Chapman \& Hall, 1993.
\end{thebibliography}

\end{document}
